\documentclass[12pt, a4paper]{article}

\usepackage[T2A]{fontenc}
\usepackage[utf8]{inputenc}
\usepackage[russian]{babel}
\usepackage{amsmath, amsfonts, amssymb, amsthm}
\usepackage[margin=2cm]{geometry}
\usepackage{parskip}

% Обозначение для мощности / количества элементов (вертикальные черты)
\newcommand{\card}[1]{\left| #1 \right|}

\begin{document}

% ==================== РАЗДЕЛ 1 ====================
\section*{Теория мер и полуколец}

$R$ --- полукольцо на $X \iff$ 
\begin{enumerate}
    \item $\varnothing \in R$
    \item $A, B \in R \implies A \cap B \in R$
    \item $A, B \in R \implies A \setminus B = \bigsqcup\limits_{i=1}^n C_i, \text{ где } C_i \in R$
\end{enumerate}

\subsection*{Теорема Каратеодори}
$X$ --- множество, $R$ --- полукольцо на $X$; если $\mu_0$ --- $\sigma$-аддитивная мера на $R$, то $\exists \mu$ --- $\sigma$-аддитивная мера на $\sigma(R)$, где $\mu(R') = \mu_0(R')$ при $R' \subseteq R$.

\subsection*{Теорема о единственности продолжения}
$\mu_1, \mu_2$ --- $\sigma$-аддитивные меры на $\sigma(S)$, где $S$ --- $\pi$-система. \\
Если:
\begin{enumerate}
    \item $\mu_1(A) = \mu_2(A)$, где $A \in S$
    \item $\exists X = \bigcup\limits_{i=1}^\infty C_i$, где $C_i \in S$ и $\mu_1(C_i) = \mu_2(C_i) < +\infty$
\end{enumerate}
Тогда $\mu_1 \equiv \mu_2$ на $\sigma(S)$.

\bigskip
\hrule
\bigskip

% ==================== ЗАДАЧА 1 ====================
\section*{Задача 1}

Пусть $F = \{ [0, 1] \} \cup \{ \{x\} \mid x \in [0, 1] \}$

$\sigma(F) = \{ A \subseteq [0, 1] \mid A \text{ счетно или косчетно} \}$

Пусть
\[
\mu_1 = 
\begin{cases}
\card{A}, & \text{если } A \text{ конечно} \\
+\infty, & \text{иначе}
\end{cases}
\quad \text{— } \sigma\text{-адд на } \sigma(F)
\]

\[
\mu_2 = 
\begin{cases}
\card{A}, & \text{если } A \text{ конечно или счетно} \\
+\infty, & \text{если } [0, 1] \setminus A \text{ конечно или счетно}
\end{cases}
\quad \text{— } \sigma\text{-адд на } \sigma(F)
\]

\begin{enumerate}
    \item $x \in F \setminus [0, 1] : \mu_1(x) = 1 = \mu_2(x)$ \\
    $x = [0, 1] : \mu_1(x) = +\infty = \mu_2(x)$
    
    \item $x \in \sigma(F) \text{ и } x\text{-счетно} \implies \mu_1(x) = +\infty \text{ и } \mu_2(x) = \aleph_0$
\end{enumerate}

\bigskip
\hrule
\bigskip

% ==================== ЗАДАЧА 2 ====================
\section*{Задача 2}

(!): $A, B \subseteq \mathbb{N} : C = A \cap B : \not\exists \lim\limits_{N \to +\infty} \frac{\card{C \cap [1, N]}}{N}$

Разобьём $\mathbb{N}$ на $I_k = [2^{k-1}; 2^k)$

Выберем $A, B \subseteq \mathbb{N}$ так:
\begin{enumerate}
    \item $k \vdots 2 : A \cap I_k = I_k, \; B \cap I_k = (\text{четные в } I_k) \implies A \cap B \cap I_k = B \cap I_k$
    \item $k \not\vdots 2 : A \cap I_k = (\text{нечет в } I_k), \; B \cap I_k = (\text{чет в } I_k), \; \cap = \varnothing$
\end{enumerate}

\[
\lim_{N \to +\infty} \frac{\card{A \cap [1, N]}}{N} = \frac{1}{2}, \quad \text{т.к. } I_{k \vdots 2} \text{ убывается } I_{k+1}
\]
\[
\lim_{N \to +\infty} \frac{\card{B \cap [1, N]}}{N} = \frac{1}{2}
\]

$\implies \limsup\limits_{N \to +\infty} = \frac{1}{2}$ и $\liminf\limits_{N \to +\infty} = 0 \implies \not\exists \lim\limits_{N \to +\infty} \frac{\card{C \cap [1, N]}}{N} \implies$

$\implies$ Этот класс не алгебра.

\bigskip
\hrule
\bigskip

% ==================== ЗАДАЧА 3 ====================
\section*{Задача 3}

\noindent а) Пусть $B_1 = A_1$, $B_n = A_n \setminus A_{n-1} \; ; \quad B_i \cap B_j = \varnothing$, $A_n = \bigcup\limits_{i=1}^n B_i$, $A = \bigcup\limits_{i=1}^\infty B_i$

\[
\mu(A) = \mu\left(\bigcup_{i=1}^\infty B_i\right) = \sum_{i=1}^\infty \mu(B_i) = \lim_{n \to \infty} \left( \sum_{i=1}^n \mu(B_i) \right) = \lim_{n \to \infty} \left( \mu\left(\bigcup_{i=1}^n B_i\right) \right) = \lim_{n \to \infty} \mu(A_n)
\]

\bigskip

\noindent Пусть $C_n = A_1 \setminus A_n \; ; \quad C_1 \subseteq C_2 \subseteq \dots$
\[
\bigcup_{n=1}^\infty C_n = \bigcup_{n=1}^\infty (A_1 \setminus A_n) = A_1 \setminus \bigcap_{n=1}^\infty A_n = A_1 \setminus A
\]

\[
\mu(A_1 \setminus A) = \lim_{n \to \infty} \mu(A_1 \setminus A_n) \quad \text{по а);}
\]

\noindent $A_n \subseteq A_1$ и $\mu(A_1) < \infty \implies \mu(A_1 \setminus A_n) = \mu(A_1) - \mu(A_n)$,

\noindent $A \subseteq A_1 \phantom{\text{и }\mu(A_1) < \infty} \implies \mu(A_1 \setminus A) = \mu(A_1) - \mu(A) \implies$

\[
\implies \mu(A_1) - \mu(A) = \lim_{n \to \infty} (\mu(A_1) - \mu(A_n))
\]
\[
\mu(A) = \lim_{n \to \infty} \mu(A_n)
\]

\bigskip

\noindent б) $B_n = A_n \setminus \left( \bigcup\limits_{i=1}^{n-1} A_i \right) \implies$
\begin{enumerate}
    \item[(1)] $B_n \subseteq A_n \implies \mu(B_n) \le \mu(A_n)$
    \item[(2)] $B_i \cap B_j = \varnothing$
    \item[(3)] $\bigcup\limits_{n=1}^\infty A_n = \bigcup\limits_{n=1}^\infty B_n$
\end{enumerate}

\[
\mu\left(\bigcup_{i=1}^\infty A_i\right) = \mu\left(\bigcup_{i=1}^\infty B_i\right) = \sum_{i=1}^\infty \mu(B_i) \overset{(1)}{\le} \sum_{i=1}^\infty \mu(A_i)
\]

\bigskip

\noindent (!): 
\[
\mu\left(\bigcup_{i=1}^n A_i\right) = \sum_{i=1}^n \mu(A_i) - \sum_{1 \le i < j \le n} \mu(A_i \cap A_j) + \sum_{1 \le i < j < k \le n} \mu(A_i \cap A_j \cap A_k) - \dots
\]

\[
I_{\bigcup\limits_{i=1}^n A_i} = 1 - I_{\bigcap\limits_{i=1}^n \overline{A_i}} = 1 - \prod_{i=1}^n (1 - I_{A_i}) = \sum_{i=1}^n I_{A_i} - \sum_{1 \le i < j \le n} I_{A_i \cap A_j} + \dots
\]

\[
\mu\left(\bigcup_{i=1}^n A_i\right) = \int I_{\bigcup\limits_{i=1}^n A_i} \, d\mu = \sum_{i=1}^n \mu(A_i) - \sum_{1 \le i < j \le n} \mu(A_i \cap A_j) + \dots
\]

\bigskip
\hrule
\bigskip

% ==================== ЗАДАЧА 4 ====================
\section*{Задача 4}

\[
\limsup_{n \to +\infty} A_n = \bigcap_{N=1}^\infty \bigcup_{n=N}^\infty A_n \; ; \quad B_N = \bigcup_{n=N}^\infty A_n
\]

\[
\limsup_{n \to +\infty} A_n \subseteq B_N
\]

\[
\mu(B_N) = \mu\left(\bigcup_{n=N}^\infty A_n\right) \le \sum_{n=N}^\infty \mu(A_n)
\]

\[
0 \le \mu\left(\limsup_{n \to +\infty} A_n\right) \le \mu(B_N) \le \sum_{n=N}^\infty \mu(A_n) \to 0 \quad \text{при } N \to +\infty
\]

\[
\implies \mu\left(\limsup_{n \to +\infty} A_n\right) = 0
\]

\bigskip
\hrule
\bigskip

% ==================== ЗАДАЧА 5 ====================
\section*{Задача 5}

\noindent а) (!): $\mathbb{Q} = \bigcap\limits_{i=1}^\infty U_i$, где $U_i$ --- открыто в $\mathbb{R}$.

\bigskip

\noindent $I = \mathbb{R} \setminus \mathbb{Q} \; ; \quad \mathbb{Q} \text{ --- счетно} \implies I = \bigcap\limits_{i=1}^\infty \left(\mathbb{R} \setminus \{q_i\}\right)$

\noindent Пусть (!) выполняется:
\[
\varnothing = \mathbb{Q} \cap I = \left(\bigcap_{i=1}^\infty U_i\right) \cap \left(\bigcap_{i=1}^\infty (\mathbb{R} \setminus \{q_i\})\right)
\]

\noindent \textbf{Теорема Бэра:} пусть $A_1, A_2, \dots$ --- плотны, тогда $\bigcap\limits_{i=1}^\infty A_i$ --- плотно.

\bigskip

\noindent $\mathbb{Q}$ --- плотно в $\mathbb{R} \; ; \quad \mathbb{Q} = \bigcap\limits_{i=1}^\infty U_i \implies \mathbb{Q} \subseteq U_i \implies \overline{\mathbb{Q}} \subseteq \overline{U_i}$
\[
\overline{\mathbb{Q}} = \mathbb{R} \implies \mathbb{R} \subseteq \overline{U_i} \implies \overline{U_i} = \mathbb{R} \implies U_i \text{ плотно в } \mathbb{R}
\]

\noindent Аналогично т.к. $I$ --- плотно в $\mathbb{R}$, то $\mathbb{R} \setminus \{q_i\}$ плотно в $\mathbb{R}$.

\bigskip

\noindent По теореме Бэра: $\left(\bigcap\limits_{i=1}^\infty U_i\right) \cap \left(\bigcap\limits_{i=1}^\infty (\mathbb{R} \setminus \{q_i\})\right)$ --- плотно в $\mathbb{R}$
\[
\implies \varnothing \text{ плотно в } \mathbb{R} \text{ --- противоречие.}
\]

\bigskip

\noindent б) $\displaystyle\bigcap_{i=1}^{\infty} A_i = \bigcup_{i=1}^{\infty} B_i$, где $A_i$ --- открыто, $B_i$ --- замкнуто.
    
Если $I \in F_\sigma$, то $\mathbb{Q} \in G_\delta$, но $\mathbb{Q} \notin G_\delta \Rightarrow I \notin F_\sigma$.

Возьмём $E = (\mathbb{Q} \cap A) \cup ((\mathbb{R} \setminus \mathbb{Q}) \cap B)$, где $A = (-\infty; 0)$, $B = [0; +\infty)$.

\begin{enumerate}
    \item[1)] Если $E \in F_\sigma$, то $E \cap B \in F_\sigma$.
    
    $E \cap B = (\mathbb{R} \setminus \mathbb{Q}) \cap B \notin F_\sigma \Rightarrow E \notin F_\sigma$.
    
    \item[2)] Если $E \in G_\delta$, то $E \cap A \in G_\delta$.
    
    $E \cap A = \mathbb{Q} \cap A \notin G_\delta \Rightarrow E \notin G_\delta$.
\end{enumerate}

\bigskip
\hrule
\bigskip

% ==================== ЗАДАЧА 6 ====================
\section*{Задача 6}

$A$ --- множество $\alpha$, для которых $\exists$ бесконечно много $(a, q)$:

\[
\left| \alpha - \frac{a}{q} \right| \le \frac{1}{q^3}
\]

$E_{a, q}$ --- множество $\alpha \in A$ при фиксированных $a, q$.

Если $\alpha \in E$, то $\alpha + N \in A$, где $N \in \mathbb{N}$:
\[
\begin{cases}
\alpha' = \alpha + N \\
a' = a + Nq
\end{cases}
\qquad 
\left( \gcd(a', q) = \gcd(a + Nq, q) = \gcd(a, q) \right)
\]
\[
\left| \alpha' - \frac{a'}{q} \right| = \left| (\alpha + N) - \frac{a + Nq}{q} \right| = \left| \alpha - \frac{a}{q} \right|
\]
Отображение из $[0, 1]$ биективно на $[N, N + 1]$.

\medskip

\[
\left| \alpha - \frac{a}{q} \right| \le \frac{1}{q^3} \implies \frac{a}{q} - \frac{1}{q^3} \le \alpha \le \frac{a}{q} + \frac{1}{q^3} 
\implies E_{a, q} = \left[ \frac{a}{q} - \frac{1}{q^3}, \frac{a}{q} + \frac{1}{q^3} \right] 
\implies \lambda(E_{a, q}) = \frac{2}{q^3}
\]

Для $\alpha \in [0, 1]$:
\[
\frac{a}{q} - \frac{1}{q^3} \le \alpha \le \frac{a}{q} + \frac{1}{q^3}; \; \alpha \in [0, 1] \implies -\frac{1}{q^2} \le a \le q + \frac{1}{q^2} \quad (q > 0)
\]
При условиях $0 \le a < q$ и $a \in \mathbb{Z}$:
\[
\begin{cases}
q = 1: & -1 \le a \le 2 \\
q \ge 2: & -0.25 \le a \le q - 0.25
\end{cases}
\]

Фиксируем только $q$: $B_q = \displaystyle\bigcup_{\gcd(a, q) = 1} E_{a, q}$.

\[
\limsup_{q \to +\infty} B_q = \bigcap_{N=1}^{\infty} \left( \bigcup_{q=N}^{\infty} B_q \right)
\]
\[
\limsup_{q \to +\infty} \left( \bigcup_{\gcd(a, q) = 1} E_{a, q} \right) = E = A
\]

Оценка меры:
\[
\sum_{q=1}^{\infty} \sum_{\substack{0 \le a < q \\ \gcd(a, q) = 1}} \lambda(E_{a, q}) 
\le \sum_{q=1}^{\infty} \sum_{a=0}^{q} \frac{2}{q^3} 
= \sum_{q=1}^{\infty} \frac{2(q+1)}{q^3} 
= 2 \sum_{q=1}^{\infty} \frac{1}{q^2} + 2 \sum_{q=1}^{\infty} \frac{1}{q^3} < \infty
\]

\[
\implies \mu(E \cap [0, 1]) = 0 \implies \mu(E) = 0
\]

\bigskip
\hrule
\bigskip

% ==================== ЗАДАЧА 7 ====================
\section*{Задача 7}

\[
K(a) = \bigcap_{n=1}^{\infty} E_n, \quad E_n \text{ --- замкнуто} \implies K(a) \text{ --- замкнуто}
\]

\[
d_n = \frac{d_{n-1} - a^n}{2} = \frac{1}{2^n} \left( 1 - \sum_{k=1}^{n} 2^{k-1} a^k \right) \to 0 \quad \text{при } n \to +\infty
\]

\[
x \in K(a) \implies x \in E_n, \; \forall n \in \mathbb{N} \implies x \in [a_n(x), b_n(x)]
\]
\[
a_n(x), b_n(x) \in E_n, \; \forall n \in \mathbb{N}
\]
\[
|x - a_n(x)| \le d_n, \quad |x - b_n(x)| \le d_n
\]
\[
d_n \to 0 \implies \forall \varepsilon > 0 : (x - \varepsilon, x + \varepsilon) \cap (K(a) \setminus \{x\}) \neq \varnothing
\]

Пусть $G$ --- удаляемые множества:
\[
\implies \lambda(G) = \sum_{n=1}^{\infty} 2^{n-1} a^n = a \sum_{n=1}^{\infty} (2a)^{n-1} = \frac{a}{1 - 2a}
\]
\[
\lambda(K(a)) = \lambda([0, 1]) - \lambda(G) = \frac{1 - 3a}{1 - 2a} > 0
\]

\end{document}